% Part: propositional-logic
% Chapter: syntax-and-semantics
% Section: formulas

\documentclass[../../../include/open-logic-section]{subfiles}

\begin{document}

\olfileid{pl}{syn}{fml}

\olsection{\usetoken{P}{formula} Proposisional}

!!^{formula}s logika proposisional kawangun saka
\emph{!!{propositional
    variable}s}\iftag{prvFalse}{\iftag{prvTrue}{,}{ lan} konstanta
  proposisional~$\lfalse$ kanggo kasalahan}{}\iftag{prvTrue}{\iftag{prvFalse}{
    lan}{} konstanta proposisional~$\ltrue$ kanggo kabeneran}{} nganggo
  \emph{konektif logis}.

\begin{enumerate}
\item Himpunan~$\PVar$ kang !!^a{denumerable} lan dumadi saka
  !!{propositional variable}s $\Obj p_0$,
  $\Obj p_1$, \dots
\tagitem{prvFalse}{Konstanta proposisional kanggo !!{falsity}~$\lfalse$.}{}
\tagitem{prvTrue}{Konstanta proposisional kanggo !!{truth}~$\ltrue$.}{}
\item Konektif logis:
  \startycommalist
  \iftag{prvNot}{\ycomma $\lnot$ (negasi)}{}%
  \iftag{prvAnd}{\ycomma $\land$ (konjungsi)}{}%
  \iftag{prvOr}{\ycomma $\lor$ (disjungsi)}{}%
  \iftag{prvIf}{\ycomma $\lif$ (!!{conditional})}{}%
  \iftag{prvIff}{\ycomma $\liff$ (!!{biconditional})}{}%
\item Tandha wacan: (, ), lan koma.
\end{enumerate}

Basa logika proposisional iki dilambangake kanthi $\Lang L_0$.

\iftag{defNot,defOr,defAnd,defIf,defIff,defTrue,defFalse,defEx,defAll}{%
Saliyane konektif primitif kang ditepungake ing ndhuwur, kita uga
nggunakake simbol-simbol kang \emph{ditetepake} iki:
\startycommalist
  \iftag{defNot}{\ycomma $\lnot$ (negasi)}{}%
  \iftag{defAnd}{\ycomma $\land$ (konjungsi)}{}%
  \iftag{defOr}{\ycomma $\lor$ (disjungsi)}{}%
  \iftag{defIf}{\ycomma $\lif$ (!!{conditional})}{}%
  \iftag{defIff}{\ycomma $\liff$ (!!{biconditional})}{}%
  \iftag{defFalse}{\ycomma $\lfalse$ (!!{falsity})}{}%
  \iftag{defTrue}{\ycomma $\ltrue$ (!!{truth})}}{}.

\begin{tagblock}{defNot,defOr,defAnd,defIf,defIff,defTrue,defFalse,defEx,defAll}
\begin{explain}
Simbol kang ditetepake dudu perangan resmi saka basa, nanging
ditepungake minangka cekakan ora resmi: simbol kasebut ndadekake kita bisa
nyekakake formula kang bakal dadi dawa banget yen mung nggunakake
simbol primitif. Iki cetha migunani. Nanging, paedahe kang luwih gedhe
yaiku bukti dadi luwih cekak. Yen sawijining simbol iku primitif, simbol
kasebut kudu ditangani kanthi kapisah ing bukti. Mula, saya
akeh simbol primitif, saya dawa bukti kita.
\end{explain}
\end{tagblock}

% Alternate symbols

\begin{intro}
Mbokmenawa pamaca wis kulina karo istilah lan simbol kang beda saka
kang digunakake ing ndhuwur. Buku teks logika (lan guru) lumrahe
nggunakake $\sim$, $\neg$, lan~!\ kanggo ``negasi'', sarta $\wedge$,
$\cdot$, lan $\&$ kanggo ``konjungsi''. Simbol kang umum digunakake
kanggo ``kondisional'' utawa ``implikasi'' yaiku $\rightarrow$,
$\Rightarrow$, lan $\supset$.
\iftag{prvIff,defIff}{Simbol kanggo ``bikondisional'', ``bi-implikasi'',
  utawa ``ekuivalensi (material)'' yaiku $\leftrightarrow$,
  $\Leftrightarrow$, lan $\equiv$.}{}
\iftag{prvFalse,defFalse}{Ing pustaka logika basa Inggris, simbol
  $\lfalse$ duwe maneka jeneng: ``falsity'' (kasalahan), ``falsum''
  (istilah Latin), ``absurdity'' (absurditas), utawa ``bottom'' (ngisor).}{}
\iftag{prvTrue,defTrue}{Ing pustaka logika basa Inggris, simbol
  $\ltrue$ duwe maneka jeneng: ``truth'' (kabeneran), ``verum'' (istilah
  Latin), utawa ``top'' (ndhuwur).}{}
\end{intro}

\begin{defn}[Formula]
\ollabel{defn:formulas}
Himpunan~$\Frm[L_0]$ saka \emph{!!{formula}s} logika proposisional
ditemtokake kanthi induktif kaya mangkene:
\begin{enumerate}
\tagitem{prvFalse}{$\lfalse$ iku !!{formula} atomik.}{}

\tagitem{prvTrue}{$\ltrue$ iku !!{formula} atomik.}{}

\item Saben !!{propositional variable}~$\Obj p_i$ iku !!{formula}
  atomik.

\tagitem{prvNot}{Yen $!A$ iku !!a{formula}, mula $\lnot !A$ iku
  !!a{formula}.}{}

\tagitem{prvAnd}{Yen $!A$ lan $!B$ iku !!{formula}s, mula $(!A \land
  !B)$ iku !!a{formula}.}{}

\tagitem{prvOr}{Yen $!A$ lan $!B$ iku !!{formula}s, mula $(!A \lor !B)$
  iku !!a{formula}.}{}

\tagitem{prvIf}{Yen $!A$ lan $!B$ iku !!{formula}s, mula $(!A \lif !B)$
  iku !!a{formula}.}{}

\tagitem{prvIff}{Yen $!A$ lan $!B$ iku !!{formula}s, mula $(!A \liff !B)$
  iku !!a{formula}.}{}

\tagitem{limitClause}{Ora ana ekspresi liyane kang kalebu !!a{formula}.}{}
\end{enumerate}
\end{defn}

\begin{explain}
Definisi !!{formula}s iki minangka \emph{definisi induktif}.
Sacara pokok, kita mbangun himpunan !!{formula}s lumantar tahap-tahap
kang cacahe tanpa wates. Ing tahap wiwitan, kita netepake yen kabeh
formula atomik iku formula; iki cocog karo sawetara kasus kapisan ing definisi,
yaiku kasus kanggo
\iftag{prvTrue}{$\ltrue$, }{}%
\iftag{prvFalse}{$\lfalse$, }{}%
$\Obj p_i$. Mula, ``!!{formula} atomik'' ateges saben !!{formula}
kanthi wujud iki.

Kasus liyane ing definisi menehi aturan kanggo mbangun !!{formula}s
anyar saka !!{formula}s kang wis kawangun. Ing tahap kapindho, kita
bisa nggunakake aturan kasebut kanggo mbangun !!{formula}s saka
!!{formula}s atomik. Ing tahap katelu, kita mbangun formula anyar saka
formula atomik lan formula kang kawangun ing tahap kapindho, banjur
sateruse. !!{formula} yaiku ekspresi apa wae kang pungkasane kawangun ing
sawijining tahap kaya mangkono, lan ora ana liyane.
\end{explain}

Nalika nulis formula $(!B \ast !C)$ kang kawangun saka $!B$, $!C$
nganggo konektif rong-papan~$\ast$, kita kerep ngilangake pasangan
tandha kurung paling njaba lan mung nulis~$!B \ast !C$.

\begin{tagblock}{defNot,defOr,defAnd,defIf,defIff,defTrue,defFalse,defEx,defAll}
\begin{defn}
Formula kang kawangun nganggo operator kang ditetepake kudu dimangerteni
kaya mangkene:

\begin{tagenumerate}{defTrue,defFalse,defNot,defOr,defAnd,defIf,defIff,defEx,defAll}

\tagitem{defTrue}{$\ltrue$ nyekakake
  \iftag{prvFalse}{$\lnot\lfalse$}{$(!A \lor \lnot !A)$ kanggo sawijining
    !!{formula} atomik~$!A$ kang tetep}.}{}

\tagitem{defFalse}{$\lfalse$ nyekakake
  \iftag{prvTrue}{$\lnot\ltrue$}{$(!A \land \lnot !A)$ kanggo sawijining
    !!{formula} atomik~$!A$ kang tetep}.}{}

\tagitem{defNot}{$\lnot !A$ nyekakake $!A \lif \lfalse$.}{}

\tagitem{defOr}{$!A \lor !B$ nyekakake
  \iftag{prvAnd}{$\lnot(\lnot !A \land \lnot !B)$}{$\lnot !A \lif
    !B$}.}{}

\tagitem{defAnd}{$!A \land !B$ nyekakake
  \iftag{prvOr}{$\lnot(\lnot !A \lor \lnot !B)$}{$\lnot (!A \lif
    \lnot !B)$}.}{}

\tagitem{defIf}{$!A \lif !B$ nyekakake
  \iftag{prvOr}{$(\lnot !A \lor !B)$}{$\lnot (!A \land \lnot !B)$}.
  Cathetan koreksi sumber OLPL-001: cabang kapisan ing sumber Inggris
  ora duwe tandha kurung buka; pasangan kurunge dibalekake.}{}

\tagitem{defIff}{$!A \liff !B$ nyekakake $(!A \lif !B) \land (!B
  \lif !A)$.}{}
\end{tagenumerate}
\end{defn}
\end{tagblock}

\begin{defn}[Identitas sintaktis]
Simbol $\ident$ mratelakake identitas sintaktis antarane rerangken
simbol, yaiku $!A \ident !B$ yen lan mung yen $!A$ lan $!B$ iku
rerangken simbol kang dawane padha lan ngemot simbol kang padha ing
saben papan.
\end{defn}

Simbol $\ident$ bisa diapit dening rerangken kang kawangun kanthi
nyambungake rerangken utawa simbol siji sawise sijine. Tuladhane,
$!A \ident (!B \lor !C)$ ateges: rerangken simbol~$!A$ padha karo
rerangken kang kawangun kanthi nyambungake, miturut urutan iki, tandha
kurung buka, rerangken $!B$, simbol $\lor$, rerangken~$!C$, lan tandha
kurung tutup. Yen mangkono, kita ngerti yen simbol kapisan $!A$ iku
tandha kurung buka, $!A$ ngemot $!B$ minangka subrerangken (wiwit saka
simbol kapindho), sawise subrerangken kasebut ana $\lor$, lan sateruse.

\end{document}
